\chapter{Green's functions}

% C1-C3 statements approved by Miguel 2026-11-26.

\begin{definition}[Retarded one-body Green's function]\label{def:green-ret}
  \lean{QM.Green.evol, QM.Green.heis, QM.Green.greenRet, QM.Green.greenZ, QM.Green.evol_apply_of_eigen,
    QM.Green.star_evol}\leanok
  \uses{def:fock}
  $G^R_{ij}(t) = -i\,\theta(t)\,\langle\{c_i(t),c_j^\dagger\}\rangle$, with $c_i(t)=e^{iHt}c_ie^{-iHt}$, in the ground state
  or a Gibbs state.
\end{definition}
% Formalized 2026-11-26 for an eigenstate of H (in particular the ground state), with e^{-iHt} the exponential of a
% bounded operator and G(z) = int_0^infty e^{izt} G^R(t) dt. The Gibbs-state version follows by linearity once the
% Gibbs state (def:gibbs, B4) is defined.

\begin{theorem}[Lehmann representation]\label{thm:lehmann}
  \lean{QM.Green.lehmannPole, QM.Green.lehmannWeight, QM.Green.greenRet_eq_sum, QM.Green.greenZ_eq_sum,
    QM.Green.lehmannWeight_posSemidef, QM.Green.sum_lehmannWeight, QM.Green.imPart, QM.Green.greenMat,
    QM.Green.neg_inv_pi_imPart_greenMat, QM.Green.tendsto_spectral, ForMathlib.poissonKernel,
    ForMathlib.tendsto_integral_poissonKernel}\leanok
  \uses{def:green-ret, thm:car}
  For $\operatorname{Im}z>0$,
  $G_{ij}(z) = \sum_{n}\frac{\langle 0|c_i|n\rangle\langle n|c_j^\dagger|0\rangle}{z-(E_n-E_0)}
             + \sum_{m}\frac{\langle 0|c_j^\dagger|m\rangle\langle m|c_i|0\rangle}{z+(E_m-E_0)}$,
  and the spectral function $A(\omega) = -\pi^{-1}\operatorname{Im}G(\omega+i0)$ is a positive matrix-valued measure of total
  mass the identity.
\end{theorem}
\begin{proof}\leanok
  Let $\varphi_n$ be an orthonormal eigenbasis of $H$ with energies $E_n$, and let $|0\rangle$ be an eigenstate with energy
  $E_0$. Then $e^{-iHt}\varphi_n=e^{-iE_nt}\varphi_n$, from the exponential series, and $\langle0|e^{-iHt}=e^{-iE_0t}\langle0|$
  because $e^{-iHt}$ has adjoint $e^{iHt}$. Inserting $\sum_n|\varphi_n\rangle\langle\varphi_n|$ gives, for $t>0$,
  $G^R_{ij}(t)=-i\sum_kW_{k,ij}\,e^{-i\omega_kt}$ with the poles $\omega_n=E_n-E_0$ and $\omega_m=E_0-E_m$. Integrating
  $\int_0^\infty e^{i(z-\omega)t}\,dt=i/(z-\omega)$ for $\operatorname{Im}z>0$ gives the formula. Each weight is a rank-one
  matrix $u\,u^\dagger$, with $u_i=\langle0|c_i|\varphi_n\rangle$ or $u_i=\langle\varphi_m|c_i|0\rangle$, so it is positive
  semidefinite. Completeness turns the sum of the weights into $\langle0|\{c_i,c_j^\dagger\}|0\rangle=\delta_{ij}$. Finally
  $-\pi^{-1}\operatorname{Im}G(\omega+i\eta)=\sum_kW_kP_\eta(\omega-\omega_k)$ with the Poisson kernel
  $P_\eta(x)=\eta/\pi(x^2+\eta^2)$. The kernel is an approximate identity (substitute $x=\eta s$ and use dominated
  convergence with $\int(1+s^2)^{-1}ds=\pi$), so $A=\sum_kW_k\,\delta_{\omega_k}$.
\end{proof}

\begin{definition}[Self-energy]\label{def:self-energy}
  \uses{thm:lehmann}
  $\Sigma(z) = G_0(z)^{-1} - G(z)^{-1}$ (Dyson equation).
\end{definition}

\begin{theorem}[Herglotz property of the self-energy]\label{thm:herglotz}
  \uses{def:self-energy}
  $\operatorname{Im}\Sigma(z)\le 0$ for $\operatorname{Im}z>0$, and $\Sigma(z)\to\Sigma_\infty$ (the Hartree–Fock self-energy)
  as $|z|\to\infty$.
\end{theorem}

\begin{theorem}[Galitskii–Migdal formula]\label{thm:galitskii-migdal}
  \uses{thm:lehmann}
  $E_0 = \tfrac12\sum_{ij}\int\! d\omega\,\theta(\mu-\omega)\,(\omega\,\delta_{ij}+h_{ij})A_{ji}(\omega)$.
\end{theorem}
